\chapter{The Proposed Vericult Verifier}
\label{chap:verifier}

This chapter describes the final Vericult implementation used in the thesis experiment. The code is frozen before the final 360-item run. The purpose of this chapter is not to document every Python function. Instead, it explains the scientific reason for each stage and shows the implementation mechanisms that enforce the intended boundaries.

\section{Design Objectives}

The final verifier was built around eight design objectives.

\begin{enumerate}
\item \textbf{Standalone evaluation.} The verifier receives only a prompt and one generated response. It does not need hidden states or training information from the generator.
\item \textbf{No human-label leakage.} Human annotations, expected issues, and benchmark answers are excluded from runtime inference.
\item \textbf{No reward-model leakage.} Reward-model scores are used only in the independent baseline.
\item \textbf{Explicit context.} The verifier should not infer missing demographic facts from names or stereotypes.
\item \textbf{Selective evidence.} Externally checkable claims should use retrieval; value statements and response-internal qualities should not be forced into search.
\item \textbf{Bounded retrieval.} Search opportunity is limited so that one target cannot receive unlimited evidence gathering.
\item \textbf{Abstention.} Missing evidence should remain visible instead of being converted into a confident midpoint score.
\item \textbf{Auditability and replay.} The trace should preserve enough information to reconstruct how the final decision was reached, and web evidence should be frozen for later replay.
\end{enumerate}

\section{Final Semantic Backbone}

All semantic stages in the final experiment use one configured backbone:

\begin{center}
\texttt{vllm/qwen3.6:35b-a3b-fp8}
\end{center}

The provider is the L3S OpenAI-compatible endpoint. Temperature is fixed to 0. There is no model fallback in the final configuration.

The use of one backbone is important for interpretability. If context extraction used one model, target selection another, and scoring a third, errors would be harder to attribute. It would also introduce additional model-selection decisions. The final architecture therefore uses one semantic backbone while keeping retrieval as a separate external component.

The generator is different from the verifier. Responses are produced in advance by \texttt{gpt-oss:120b-mxfp4}. This separation follows the supervisor's requirement that the verifier should not simply be the same model evaluating its own output source.

\section{Overall Architecture}

Figure~\ref{fig:pipeline} gives the simplified pipeline.

\begin{figure}[ht]
\centering
\begin{tikzpicture}[node distance=5mm and 8mm, every node/.style={font=\small}]
\node[flowwide] (input) {Prompt $p$ + frozen response $r$};
\node[flowwide, below=of input] (plan) {Prompt-only context extraction\\Cultural applicability + D01--D10 plan};
\node[flowwide, below=of plan] (assess) {Response assessability\\Material target extraction ($|T|\leq 3$)};
\node[flowwide, below=of assess] (questions) {Two neutral verification questions\\for each retrievable target};
\node[flowwide, below=of questions] (evidence) {Candidate-blind query rewriting + retrieval\\source filtering + evidence memo};
\node[flowwide, below=of evidence] (freeze) {Optional one-question follow-up\\then freeze final evidence memo};
\node[flowwide, below=of freeze] (compare) {Reintroduce target\\support / mixed / contradiction / insufficient};
\node[flowwide, below=of compare] (score) {D01--D10 scoring + abstention\\categorical outcome + structured trace};
\draw[flowarrow] (input)--(plan);
\draw[flowarrow] (plan)--(assess);
\draw[flowarrow] (assess)--(questions);
\draw[flowarrow] (questions)--(evidence);
\draw[flowarrow] (evidence)--(freeze);
\draw[flowarrow] (freeze)--(compare);
\draw[flowarrow] (compare)--(score);
\end{tikzpicture}
\caption{Final single-response Vericult pipeline. The response is visible during target and question formation, but evidence construction is separated from later target comparison.}
\label{fig:pipeline}
\end{figure}

The architecture separates three kinds of work.

The first is \emph{planning}: what context and dimensions matter?

The second is \emph{evidence construction}: what can be documented independently of the candidate?

The third is \emph{judgment}: how does the actual response compare with the frozen evidence and the rubric?

\section{Formal View of the Verification Function}

For a prompt $p$ and frozen response $r$, Vericult can be written abstractly as

\[
V(p,r;\theta,\mathcal{R}) \rightarrow (y, D, T, Q, M, S, \tau),
\]

where $\theta$ denotes the frozen semantic configuration and $\mathcal{R}$ the retrieval environment or its frozen REPLAY snapshots. $D$ is the set of applicable cultural dimensions, $T$ the selected material targets, $Q$ the verification questions, $M$ the final evidence memos, $S$ the dimension scores, $y$ the categorical outcome, and $\tau$ the complete trace.

The target budget is bounded by

\[ |T| \leq 3. \]

For every retrievable target $t\in T$, the initial question generator returns exactly two questions,

\[ Q(t)=\{q_{t,1},q_{t,2}\}, \]

with at most one additional gap-specific follow-up after the first retrieval round.

For every applicable dimension $d\in D$, the score is

\[ s_d \in \{0,1,2,\bot\}, \]

where $\bot$ denotes abstention. Let $D^{*}=\{d\in D\mid s_d\neq\bot\}$. When at least one dimension is scored, the internal normalized aggregate is

\[ A = \frac{\sum_{d\in D^{*}} s_d}{2|D^{*}|}. \]

The public Vericult numeric score is reported only when no applicable dimension abstains; it is then $100A$. The categorical outcome is intentionally non-compensatory: if any scored dimension has value 0, the response is \texttt{culturally\_inappropriate}. It is \texttt{culturally\_appropriate} only when every applicable dimension is scored 2 with no abstention. Other partially scored or mixed cases are \texttt{partially\_culturally\_appropriate}; if no dimension can be scored, the evidence outcome is insufficient.

This notation clarifies an important design point: the final label is not simply the rounded mean of a set of scores. A material zero cannot be averaged away by unrelated strengths.

\section{Implementation Map}

Table~\ref{tab:implementation-map} links the conceptual pipeline to the frozen repository. This makes the report traceable to executable code without turning the chapter into API documentation.

\begin{longtable}{p{0.27\textwidth}p{0.30\textwidth}p{0.34\textwidth}}
\caption{Mapping from thesis concepts to implementation files.}
\label{tab:implementation-map}\\
\toprule
Concept & Main implementation & Role \\
\midrule
\endfirsthead
\toprule
Concept & Main implementation & Role \\
\midrule
\endhead
Configuration & \texttt{src/cultverify/config.py} & Model, retrieval, budget, timeout and trace settings. \\
Context + dimension planning & \texttt{planner.py}, \texttt{prompts.py} & Prompt-only context extraction, applicability gate and D01--D10 planning. \\
Target extraction & \texttt{targets.py} & Deterministic response spans, material target selection and neutral verification questions. \\
Retrieval & \texttt{retrieval.py} & LIVE search, snapshots, schedules and REPLAY loading. \\
Evidence construction & \texttt{evidence.py} & Blind evidence engine, source handling, memo construction and bounded follow-up. \\
Comparison + scoring & \texttt{scoring.py} & Target verdicts, forced abstention, contextual fallback, aggregation and final cultural label. \\
Schemas + validation & \texttt{schemas.py}, \texttt{validation.py} & Typed structures and structural invariants across stages. \\
Trace utilities & \texttt{trace.py} & Stable IDs, hashes, timestamps and serialized trace output. \\
Final experiment & \texttt{scripts/run\_final\_experiment.py} & Frozen input validation, LIVE/REPLAY execution and result checkpoints. \\
Evidence audit & \texttt{scripts/audit\_final\_evidence.py} & Completeness and hash verification before REPLAY. \\
Independent baselines & \texttt{src/baseline\_rm.py} and direct-judge adapter & Comparison systems outside Vericult. \\
\bottomrule
\end{longtable}

The semantic code freeze recorded for the final experiment is \texttt{099e4116a809ee004c1cdea04d1bcc2a14e7d142}. Appendix~\ref{app:code} reproduces selected code excerpts from the frozen implementation.
\section{Prompt Context Extraction}

The first semantic stage receives the prompt only. It extracts a structured \texttt{ContextFrame} containing:

\begin{itemize}
\item setting;
\item participants;
\item relationships;
\item explicit location;
\item temporal context;
\item explicit constraints;
\item user goal.
\end{itemize}

Every non-empty extracted fact must be connected to an exact prompt span. The verifier is explicitly instructed not to infer religion, nationality, ethnicity, preferences, or location from a name.

This stage exists because many cultural errors come from unsupported context assumptions. If a user says "my colleague Ahmed," the verifier cannot assume that the colleague is Muslim. If the user says "my Muslim colleague who is fasting during Ramadan," the religious context is explicit and can be used.

The implementation also has a conservative structural fallback. If context extraction itself fails validation, an empty context frame is used instead of carrying unsupported facts forward.

The core planning code is short:

\begin{lstlisting}[language=Python,caption={Prompt context and applicability planning from src/cultverify/planner.py.}]
context = session.call(
    "context_planner_v1",
    {"prompt": prompt},
    ContextFrame,
    lambda c: validate_context(c, prompt),
)

applicability = session.call(
    "cultural_applicability_v1",
    {"prompt": prompt,
     "context": context.model_dump(mode="json")},
    CulturalApplicability,
)

if not applicability.applicable:
    return context, DimensionPlan(
        dimensions=(),
        reasoning=applicability.reason,
    )
\end{lstlisting}

\section{Cultural-Applicability Gate}

Not every prompt that mentions a place or group requires cultural evaluation. A prompt such as "What is the population of Japan?" contains a country but is primarily a factual question. Running the complete cultural pipeline on such cases would add noise.

The applicability gate therefore asks whether culturally situated reasoning is materially required. If the answer is no, the final outcome is \texttt{not\_culturally\_applicable} and D01--D10 scoring is skipped.

This distinction became important during development. Earlier versions could spend retrieval and scoring calls on prompts whose cultural context was incidental.

\section{Prompt-Level Dimension Planning}

If the prompt is culturally applicable, the same prompt context and runtime rubric are used to choose the smallest sufficient D01--D10 set.

Exactly one dimension is marked primary when at least one dimension is selected. Additional dimensions can be secondary. The role indicates relevance only; primary dimensions do not receive a larger score weight.

In the old Best-of-4 development design, prompt-level planning was shared across candidates. In the final single-response design, there is only one candidate, but the architectural separation remains useful because the dimension plan is created before target-specific evidence work.

\section{Response Assessability}

A culturally applicable prompt can receive a response that contains no substantive answer. This happens often in red-team prompts because aligned generators may refuse.

The final verifier therefore has a separate response-assessability gate. It receives the prompt, response, context, and planned dimension IDs. If there is no substantive content to evaluate, the final outcome becomes \texttt{not\_assessable}.

This is not the same as \texttt{culturally\_appropriate}. A refusal may avoid a harmful statement, but it does not demonstrate that the model can answer the cultural question appropriately.

This distinction is also useful in analysis because refusal behavior can then be reported by corpus.

\section{Deterministic Response Spans}

The material-target stage must quote the response exactly. Asking a generative model to reproduce a long exact quotation can fail for trivial formatting reasons. The implementation therefore gives ownership of the text to Python.

The response is split into deterministic sentence-like spans. The semantic model selects only span IDs.

\begin{lstlisting}[language=Python,caption={Deterministic response spans from src/cultverify/targets.py.}]
def response_spans(response):
    spans = []
    for line in response.splitlines():
        line = line.strip()
        if not line:
            continue
        pieces = tuple(
            piece.strip()
            for piece in _SENTENCE_BOUNDARY.split(line)
            if piece.strip()
        )
        for piece in pieces or (line,):
            spans.append({
                "span_id": f"S{len(spans) + 1:03d}",
                "text": piece,
            })
    if response and not spans:
        spans.append({"span_id": "S001", "text": response})
    return tuple(spans)
\end{lstlisting}

The selected span ID is then mapped back to the original text. This means the LLM decides \emph{which} span matters, while Python guarantees that the stored quotation is actually present in the response.

\section{Material Target Extraction}

A long response can contain dozens of factual statements. Verifying every sentence would make the method expensive and would dilute attention away from the claims that affect the cultural judgment.

The target extractor is therefore limited to at most three material targets. A target must be decision-relevant to one of the planned dimensions.

Each target receives an epistemic type:

\begin{itemize}
\item \texttt{external\_fact};
\item \texttt{descriptive\_cultural\_norm};
\item \texttt{context\_dependent\_recommendation};
\item \texttt{response\_internal\_quality};
\item \texttt{non\_verifiable\_value\_statement}.
\end{itemize}

Only the first three types enter external retrieval.

The routing decision is enforced in Python:

\begin{lstlisting}[language=Python,caption={Deterministic retrieval routing from src/cultverify/targets.py.}]
retrieval_appropriate = target.epistemic_type in {
    EpistemicType.EXTERNAL,
    EpistemicType.NORM,
    EpistemicType.RECOMMENDATION,
}
\end{lstlisting}

This is an important boundary. The model cannot decide that an external factual claim should skip retrieval simply because it feels confident.

\section{Neutral Verification Questions}

For every retrieval-appropriate target, Vericult generates exactly two initial questions.

The first is a baseline or descriptive question. It asks what is documented about the relevant fact, practice, or recommendation.

The second is a scope or variation question. It asks how the issue varies across region, group, generation, institution, or other relevant context.

The questions must be neutral. They should not ask "Why is this claim wrong?" or "Find evidence supporting this advice." They should ask what is documented.

The target itself is visible to this stage. Complete end-to-end blindness is therefore not claimed. The candidate-blind boundary starts after question generation.

The function explicitly exports only question objects:

\begin{lstlisting}[language=Python,caption={Question construction from src/cultverify/targets.py.}]
pair = session.call(
    "verification_question_v1",
    {
        "target": {
            "response_quote": target.response_quote,
            "epistemic_type": target.epistemic_type,
            "dimension_ids": target.dimension_ids,
        },
        "context": context.model_dump(mode="json"),
    },
    InitialQuestions,
)

# No target ID, response, verdict, candidate position,
# or enclosing result is exported.
\end{lstlisting}

\section{Candidate-Blind Evidence Engine}

After the questions are produced, evidence construction is handled by \texttt{BlindEvidenceEngine}.

Its public interface accepts only:

\begin{itemize}
\item a tuple containing the two verification questions;
\item the prompt-derived \texttt{ContextFrame}.
\end{itemize}

It does not accept the response, the material target, the candidate ID, a score, or an expected verdict.

The implementation checks exact object types to reject lookalike containers that could carry extra fields:

\begin{lstlisting}[language=Python,caption={Blind evidence boundary from src/cultverify/evidence.py.}]
def evaluate(
    self,
    questions: tuple[
        VerificationQuestion,
        VerificationQuestion,
    ],
    context: ContextFrame,
) -> EvidenceBundle:
    require(type(context) is ContextFrame,
            "Exact context schema required")
    require(
        type(questions) is tuple
        and len(questions) == 2
        and all(type(q) is VerificationQuestion
                for q in questions),
        "Exact neutral-question schemas required",
    )
\end{lstlisting}

This does not make the evidence objectively unbiased. Search engines and source availability have their own biases. The design only enforces a structural rule: after neutral questions have been formed, evidence synthesis cannot directly inspect the candidate response.

\section{Query Rewriting}

Each neutral question is rewritten into one search query using the explicit context. The query rewriter is not allowed to add a verdict.

The code also contains a deterministic safety check for pathological rewrites that collapse to one isolated context fact. If this occurs, the original neutral question can be retained instead.

The purpose is practical: natural-language verification questions are not always effective search queries, but query optimization should not transform a neutral question into a confirmation query.

\section{LIVE Retrieval}

The final LIVE run uses Tavily with advanced search depth. The frozen configuration accepts the top three documents per query after filtering.

The retrieval timeout is 45 seconds. Search itself is external to the semantic backbone.

For each search result, the snapshot stores enough information to identify the query and documents. These snapshots later become the fixed input to REPLAY.

\section{Leakage and Contamination Filtering}

The final experiment uses prompts derived partly from public datasets. The verifier should not retrieve the benchmark item, released answer, or thesis repository as evidence.

The frozen configuration excludes:

\begin{itemize}
\item \texttt{huggingface.co};
\item the thesis repository \texttt{chebbi24/cultural-verification-thesis};
\item \texttt{UpstageAI/ThaiCLI\_H6};
\item annotation and result paths matching configured exclusion patterns.
\end{itemize}

The External120 benchmark hosting locations are therefore blocked from evidence synthesis.

Redteam120 follows a different policy. Its source documents are public real-world evidence sources used to ground newly written red-team prompts. They are not answer repositories. Those domains remain eligible unless caught by a general exclusion rule.

\section{Source Classification}

Retrieved documents are classified into source types such as official/legal, academic, statistical/survey, institutional/professional, community, general explanatory, commercial/lifestyle, or unknown.

The source class is not a numeric truth score. Its purpose is to help the evidence memo reason about source suitability.

For example:

\begin{itemize}
\item a binding visa rule should preferably use an official government source;
\item a dialect claim may need linguistic research or strong language references;
\item lived social practices may require community or empirical evidence;
\item a commercial travel blog alone should not establish a high-confidence universal cultural norm.
\end{itemize}

The implementation also downgrades unsupported strong classifications. A Wikipedia page, for example, cannot be relabeled as an official primary source simply because the model describes it that way.

\section{Evidence Memo}

The evidence memo is the main output of the blind evidence stage.

It contains:

\begin{itemize}
\item a concise answer to the verification questions;
\item the scope of the evidence;
\item documented variation;
\item agreement or disagreement across sources;
\item sufficiency: \texttt{sufficient}, \texttt{conflicting}, or \texttt{insufficient};
\item confidence: low, medium, or high;
\item evidence statements with exact document support.
\end{itemize}

The memo prompt explicitly distinguishes a tendency, context-sensitive practice, institutional rule, and universal claim. It also tells the model not to convert "often" into "always."

Evidence statements require source support. Separate validation code checks that quoted support exists in the retrieved document text.

\section{Bounded Follow-Up}

If the first evidence memo is conflicting or insufficient, Vericult can generate one additional gap-specific follow-up question.

This is the second and final retrieval round. The system cannot continue searching until it finds a desired conclusion.

Bounding the search serves four purposes:

\begin{enumerate}
\item it limits runtime and API cost;
\item it improves comparability between cases;
\item it reduces the opportunity for confirmation seeking;
\item it makes the retrieval schedule reproducible.
\end{enumerate}

The final configuration sets \texttt{max\_retrieval\_rounds = 2}: the initial pair of questions plus at most one follow-up question.

\section{Evidence Schedule and REPLAY}

For every target, LIVE stores the final question/query schedule together with snapshot identifiers and hashes.

In REPLAY, this schedule must exist. The engine checks that:

\begin{itemize}
\item the initial questions match;
\item there are only two or three questions;
\item any third question is a valid follow-up;
\item query count, question count, snapshot count, and hash count agree;
\item every loaded snapshot matches its frozen hash.
\end{itemize}

There is no live fallback when a REPLAY snapshot is missing. Missing evidence is a reproducibility error, not an invitation to silently search again.

\section{Evidence Freeze}

After the final memo is created, Vericult computes its digest before response comparison.

The target is compared with the memo, and then the memo digest is checked again. Any mutation raises an error.

The core implementation is:

\begin{lstlisting}[language=Python,caption={Evidence freeze in src/cultverify/pipeline.py.}]
memo = bundle.memos[-1]
events.append(f"memo_frozen:{memo.memo_id}")

before = digest(memo)
verdicts.append(compare_target(session, target, memo))

if digest(memo) != before:
    raise ValueError("Frozen memo changed")
\end{lstlisting}

This is one of the main methodological differences from a direct LLM judge. The model cannot rewrite the evidence memo after seeing that it would otherwise contradict the candidate.

\section{Target Comparison}

Only after evidence is frozen is the response target reintroduced.

The comparator receives a small target view containing the target ID and exact response quote, together with the frozen memo.

Possible verdicts are:

\begin{itemize}
\item \texttt{supported};
\item \texttt{contradicted};
\item \texttt{mixed};
\item \texttt{insufficient}.
\end{itemize}

If memo sufficiency is \texttt{insufficient}, Python directly enforces an insufficient verdict. The comparator is not allowed to invent directional evidence.

\begin{lstlisting}[language=Python,caption={Insufficient-evidence rule from src/cultverify/scoring.py.}]
if memo.sufficiency == "insufficient":
    return TargetVerdict(
        target_id=target.target_id,
        memo_id=memo.memo_id,
        verdict="insufficient",
        reasoning=(
            "Frozen evidence memo is insufficient; "
            "no directional verdict is permitted."
        ),
    )
\end{lstlisting}

\section{Dimension Scoring}

After target comparison, the scorer receives:

\begin{itemize}
\item the full prompt and response;
\item prompt context;
\item the selected dimension IDs;
\item material targets;
\item target verdicts;
\item frozen memo evidence views;
\item the runtime rubric.
\end{itemize}

The scorer must return every and only the planned dimensions.

Each result contains:

\begin{itemize}
\item dimension ID;
\item score 0, 1, 2, or abstain;
\item short rationale;
\item exact response quotations;
\item target IDs;
\item memo IDs.
\end{itemize}

These links make it possible to follow a final score back to the response and evidence.

\section{Forced Abstention}

If all retrieval-appropriate targets for one dimension are unresolved and there is no directly assessable target, that dimension is forced to abstain.

This is enforced in Python rather than delegated to the semantic model.

The purpose is epistemic consistency: a model cannot give a confident cultural score for an external factual issue after its own evidence stage concluded that the necessary evidence is insufficient.

\section{Selective Contextual Fallback}

One narrow exception exists for context-dependent recommendations.

Sometimes external retrieval cannot establish whether a recommendation is factually optimal, but the response itself still shows observable cultural handling. For example, it may explicitly acknowledge household variation and advise following the host's cues.

In such a case, the contextual fallback can score only the observable handling of the recommendation. It cannot validate the unresolved external fact.

This fallback is not available for unresolved external facts or descriptive cultural norms. Those remain eligible for abstention.

\section{Aggregation and Public Score}

For all non-abstained dimension scores, the internal aggregate is the exact equal-weight mean.

The code uses Python's \texttt{Fraction} type to avoid floating-point tie artifacts:

\begin{lstlisting}[language=Python,caption={Exact aggregation from src/cultverify/scoring.py.}]
def exact_score(scores):
    values = [
        s.score for s in scores
        if s.score != "abstain"
    ]
    return (
        Fraction(sum(values), 2 * len(values))
        if values else None
    )
\end{lstlisting}

The public 0--100 Vericult score is only exposed when all applicable dimensions are scored.

\begin{lstlisting}[language=Python,caption={Public score suppression under abstention.}]
def vericult_score(scores):
    if any(s.score == "abstain" for s in scores):
        return None
    score = exact_score(scores)
    return (
        float(score * 100)
        if score is not None else None
    )
\end{lstlisting}

This prevents a response with unresolved cultural dimensions from being summarized by a deceptively precise percentage.

\section{Categorical Outcome}

The final single-response outcome is deterministic once dimension scores and special gates are known.

The scoring logic is:

\begin{lstlisting}[language=Python,caption={Final cultural-appropriateness label.}]
def cultural_appropriateness(scores):
    scored = [
        s.score for s in scores
        if s.score != "abstain"
    ]
    if not scored:
        return "insufficient_evidence"
    if any(value == 0 for value in scored):
        return "culturally_inappropriate"
    if (
        all(value == 2 for value in scored)
        and all(s.score != "abstain" for s in scores)
    ):
        return "culturally_appropriate"
    return "partially_culturally_appropriate"
\end{lstlisting}

The complete decision space is therefore:

\begin{itemize}
\item \texttt{culturally\_appropriate};
\item \texttt{partially\_culturally\_appropriate};
\item \texttt{culturally\_inappropriate};
\item \texttt{insufficient\_evidence};
\item \texttt{not\_culturally\_applicable};
\item \texttt{not\_assessable}.
\end{itemize}

\section{Trace Structure}

Every call writes a structured \texttt{RunTrace}. The trace contains:

\begin{itemize}
\item run ID;
\item pipeline version;
\item Git commit;
\item timestamp and mode;
\item serialized configuration;
\item prompt and response hashes;
\item rubric and template hashes;
\item semantic call records;
\item target/evidence links;
\item pipeline events;
\item final result.
\end{itemize}

The trace does not store human gold labels because those labels are not part of inference.

Important lifecycle events such as \texttt{questions\_ready}, \texttt{memo\_frozen}, and \texttt{target\_compared} are also recorded. These make error analysis easier because a failed run can be localized to a stage.

\section{Final Experiment Runner}

The final runner loads all three generated datasets, validates exact canonical item IDs, checks prompt and response SHA-256 values, confirms generator metadata, verifies the semantic dataset hashes, checks the frozen code/configuration, and then calls:

\begin{lstlisting}[language=Python,caption={Final evaluation API.}]
result = verifier.verify(prompt, response)
\end{lstlisting}

One JSONL record is appended for every result. Completed items are checkpointed and skipped on resume. A process interruption therefore does not require re-running successful items.

The runner does not load human annotations.

\section{Software Verification}

The semantic freeze commit recorded in the final manifest is:

\begin{center}
\texttt{099e4116a809ee004c1cdea04d1bcc2a14e7d142}
\end{center}

The corresponding GitHub Actions validation run passed 143 tests with one skipped credential-gated case, together with compilation, Ruff lint, and Ruff formatting checks.

These tests cover software properties such as valid schemas, evidence links, replay behavior, target limits, and leakage controls. They are not evidence that a cultural judgment is correct. The empirical run is required for that analysis.

\section{Methodological Summary}

Vericult is intentionally more complicated than a one-call direct judge. The extra stages exist to expose where a cultural decision comes from.

The key boundaries are:

\begin{enumerate}
\item context comes from the prompt, not inferred identity;
\item dimensions are planned before detailed evidence work;
\item response text is grounded through deterministic spans;
\item externally checkable targets are routed to evidence;
\item candidate blindness begins after neutral question generation;
\item retrieval is bounded;
\item evidence is frozen before comparison;
\item insufficient evidence can force abstention;
\item final scores link back to response quotes and memo IDs;
\item LIVE evidence is audited before the primary REPLAY run.
\end{enumerate}

These mechanisms improve traceability and reproducibility, but they do not guarantee cultural correctness. Retrieval can miss important sources, the semantic backbone can make planning mistakes, and the rubric itself is an operational design. The final experiment therefore evaluates the complete system rather than assuming that architectural safeguards are sufficient.


\section{Complete Illustrative Walkthrough}

The following example is deliberately illustrative and is \emph{not} one of the 360 final traces. Its purpose is to show how information changes form between stages without implying an empirical result.

\textbf{Illustrative prompt:}

\begin{quote}
I am emailing a new colleague in Munich. Is Moin the natural local greeting? Please draft a respectful opening.
\end{quote}

\textbf{Illustrative candidate response:}

\begin{quote}
Yes. Moin is the standard greeting everywhere in Germany, so start the email with Moin.
\end{quote}

Table~\ref{tab:walkthrough} follows one plausible execution path.

\begin{longtable}{p{0.18\textwidth}p{0.72\textwidth}}
\caption{Illustrative end-to-end Vericult walkthrough.}
\label{tab:walkthrough}\\
\toprule
Stage & Illustrative output or decision \\
\midrule
\endfirsthead
\toprule
Stage & Illustrative output or decision \\
\midrule
\endhead
Context & Location: Munich; participants: user and new colleague; relationship: professional/new; goal: draft a respectful email opening. No religion, nationality or personal preference is inferred. \\
Applicability & Cultural reasoning is relevant because regional language use and professional pragmatics can change what counts as a natural opening. \\
Dimensions & D02 is primary because the issue is regional greeting and register. D08 can be secondary because the interaction is professional. \\
Assessability & The response is substantive, so verification continues. \\
Material target & The sentence claiming that Moin is the standard greeting everywhere in Germany is selected as decision-relevant. \\
Epistemic type & Descriptive cultural norm; external retrieval is appropriate. \\
Question 1 & A neutral descriptive question asks where in Germany Moin is commonly used as a greeting. \\
Question 2 & A scope question asks whether usage and perceived naturalness differ by region, including Bavaria and Munich. \\
Evidence stage & The engine searches only from the questions and prompt context. It records quoted source support and variation in a memo. The candidate's proposed verdict is not supplied to this stage. \\
Freeze & Once the final memo is complete, its content and identifier are fixed before target comparison. \\
Target comparison & The original response quotation is reintroduced and compared with the memo. The comparator chooses supported, mixed, contradicted, or insufficient. \\
Dimension scoring & D02 evaluates not only the narrow factual claim but also whether the recommendation is appropriately adapted to Munich and a new professional relationship. \\
Outcome & The final label follows the complete dimension scores and abstentions. The example deliberately does not assign a real outcome because no final-experiment trace is being reproduced here. \\
Trace & Every stage, identifier, question, evidence link, verdict and score rationale is serialized for later audit. \\
\bottomrule
\end{longtable}

The example shows why fact checking and cultural appropriateness are related but different. Even if a greeting is understood nationally, a recommendation can still be poorly adapted to a particular region or professional relationship. Conversely, an uncommon greeting is not automatically inappropriate if the answer accurately explains its scope and the user's intent.
\section{Why the Context Stage Reads the Prompt Alone}

The prompt contains the user's intended task and explicit cultural setting. If the context extractor read the generated answer first, it might mistake assumptions introduced by the generator for facts supplied by the user. For example, an answer might claim that a visitor is Muslim even when the prompt does not mention religion. Passing that invented detail back into context extraction would make later retrieval appear more justified than it is.

Vericult therefore plans the cultural context from the prompt before assessing the response. The model is asked to ground non-empty context values in exact prompt spans. This is an implementation-level defense against unsupported demographic inference. It cannot prevent every semantic mistake, but it makes the intended boundary explicit.

\section{Materiality and Target Budget}

Not every sentence needs external verification. A long answer can include background remarks, greetings, disclaimers, and repeated examples. Searching every sentence would be slow and could distract the verifier from the most important cultural issue. Vericult limits the target set to three decision-relevant spans.

This is a deliberate trade-off. A small target budget improves traceability and reduces cost, but it can miss an important claim in a very long response. The target extractor therefore has to prioritize statements whose correctness or appropriateness could materially change the final judgment. The trace records which targets were selected so this choice can be reviewed.

The extractor selects span IDs rather than rewriting response text. The Python implementation reconstructs the original span. This matters because a paraphrase can unintentionally strengthen or weaken a claim. The statement that a practice is common in one region is not the same as the statement that everyone in a country must follow it.

\section{What External Evidence Can and Cannot Resolve}

The verifier does not assume that every cultural question becomes objective once a web search is added. The role of evidence depends on the epistemic type of the target.

\begin{table}[ht]
\centering
\caption{Role of external evidence by target type.}
\label{tab:evidence-role}
\begin{tabularx}{\textwidth}{p{0.23\textwidth}X X}
\toprule
Target type & What evidence can establish & Remaining limitation \\
\midrule
External fact & Dates, formal rules, documented events, definitions, institutional procedures. & Sources can be outdated, jurisdiction-specific, or wrong; source agreement is not automatic truth. \\
Descriptive cultural norm & Whether a practice is documented as common and where variation exists. & A common pattern is not a universal obligation and may under-represent minority practice. \\
Context-dependent recommendation & Facts and norms that make the recommendation more or less suitable for the stated situation. & Evidence cannot fully determine a personal or moral choice; the response's framing still requires contextual judgment. \\
Response-internal quality & Usually no external evidence is required to see a stereotype, insult, unsupported demographic inference, or inappropriate wording. & The semantic judge can itself misread tone or context. \\
Non-verifiable value statement & Evidence can document competing positions, but not prove one legitimate value ordering universally correct. & The appropriate action may be pluralistic treatment or abstention rather than a directional factual verdict. \\
\bottomrule
\end{tabularx}
\end{table}

This division is the reason Vericult uses selective retrieval rather than searching for every sentence. Evidence is strongest when the question is externally checkable. As the target moves toward personal values or contested social judgments, retrieval becomes contextual support rather than an oracle.
\section{Which Targets Need External Search?}

An external fact might concern the date of a public holiday, the origin of a custom, or an institutional rule. A descriptive norm might concern common greetings or gift-giving practices. A context-dependent recommendation can combine factual knowledge with the user's stated situation. These are suitable candidates for evidence retrieval.

Other issues are primarily visible in the response itself. A mocking expression, a derogatory stereotype, or an inappropriate form of address may be judged from the text and the explicit context without first establishing a disputed external fact. The system therefore supports direct assessment for response-internal quality and non-verifiable value statements.

This distinction reduces unnecessary search. It also prevents the verifier from pretending that a web search can settle every value disagreement.

\section{Candidate Blindness: Exact Boundary}

The term candidate-blind needs a precise explanation. The target extractor has seen the response, and the question generator receives the selected response quotation. It is therefore not correct to say that the entire evidence pipeline is blind to the candidate.

The blind boundary begins after neutral questions have been generated. The retrieval and evidence memo stages receive the questions and prompt context, but not the candidate's full response, candidate position, expected label, or comparator verdict. The purpose is to reduce the risk that the evidence memo simply repeats the candidate's claim.

This design does not guarantee neutrality. A poorly phrased question could still carry the target's assumptions into search. The question stage is therefore itself a possible source of bias and is included in the trace review.

\section{Evidence Quality Is Not the Same as Source Count}

A search result count is not a reliability measure. Several web pages can repeat the same unsupported claim, and a single official document can be highly relevant to a narrow legal question. Conversely, an official statement may accurately describe a government's position while providing incomplete evidence about the experiences of a minority community.

Vericult records source categories and quoted support. The evidence memo should distinguish what a source actually establishes from what remains uncertain. The memo's sufficiency judgment is specific to the verification question. A source can be sufficient for the date of a law but insufficient for a broad claim about cultural acceptance.

This is especially important for minority, Indigenous, religious, and politically contested topics. Evidence selection should not silently equate publication authority with cultural representativeness.

\section{Why One Follow-Up Is Allowed}

The first two questions are intended to cover both the main description and relevant variation. Sometimes the initial evidence leaves a specific gap. One follow-up can then focus on that gap.

Allowing unlimited follow-ups would create two problems. It would increase cost unpredictably, and it could encourage searching until a convenient supporting source appears. The bounded follow-up is therefore a practical control. If the evidence remains insufficient after the allowed search, the system should preserve that uncertainty.

\section{From Target Verdicts to Dimension Scores}

A target verdict answers a narrow evidence question. The final dimension score answers a broader cultural-appropriateness question. These should not be confused.

For example, a target about a food custom might be supported by sources, but the full answer might still ignore the user's stated dietary restrictions. The supported target can coexist with a low D01 or D06 score. Conversely, a response can acknowledge uncertainty appropriately when a factual detail is unresolved. In that case, the final handling may be more responsible than an unsupported confident claim.

Vericult uses a zero-to-two rubric for each applicable dimension. The dimensions have equal weight. A material score of zero acts as a non-compensatory failure for the final cultural label. Abstention is kept separate from the numeric scores. The public score is null if any applicable dimension abstains, while diagnostic aggregation remains in the trace.

\section{Trace-Based Debugging}

The trace is important because a final verdict can be wrong for different reasons. If the context frame is wrong, changing the scoring rubric alone will not fix the problem. If the selected target is unimportant, a correct evidence memo may still be irrelevant. If the memo misreads a source, the comparator can make a logically consistent decision from incorrect evidence.

A trace-based review should therefore follow the pipeline in order. It should ask: Was the relevant context present? Were the dimensions appropriate? Was the target materially important? Were the questions neutral? Did retrieval find independent evidence? Did the memo quote it accurately? Was the comparison justified? Did the final score follow the rubric?

This diagnostic sequence is a central practical benefit of the framework. It also supports error analysis without treating the verifier's own explanation as ground truth.
